\chapter{Introduction}
This chapter provides an introduction to the project work. It presents the context of global logistics chokepoint vulnerabilities and the impact of shipping friction on Indian trade. The chapter details the literature survey conducted in this domain, highlights the motivation and problem statement, and outlines the objectives of the work. Finally, it outlines the brief methodology, project constraints, and the organization of the report.

\section[Introduction]{\textbf{Introduction}}
Global trade is heavily dependent on maritime transport, which accounts for over 80\% of international trade volume. This network of global commerce relies on critical shipping corridors and narrow passages known as maritime chokepoints, such as the Suez Canal, the Bab-el-Mandeb Strait (Red Sea), and the Strait of Hormuz. Because these corridors carry a high density of cargo, any local disruption—whether due to geopolitical conflict, piracy, labor disputes, or weather gales—can introduce significant logistics friction. This friction manifests as transit delays, increased insurance premiums, and higher freight costs, which cascade down the supply chain to final consumption markets.

Indian export-import (EXIM) activities are particularly vulnerable to these disruptions. A major portion of India's imports (such as oil and industrial raw materials) and exports (such as engineering goods and textiles) pass through the Red Sea and the Persian Gulf. Disruptions in these regions force shipping lines to take circuitous bypass routes (such as around the Cape of Good Hope), adding thousands of nautical miles and weeks to transit schedules. Traditional static supply chain management systems are ill-equipped to handle these dynamic events. This project builds a five-layer Agentic Digital Twin framework that uses graph theory and LLM-based agentic intelligence to autonomously detect logistics friction, simulate shock propagation, and optimize alternative paths to protect local Indian Small and Medium Enterprise (SME) markets.

\section[Literature Review]{\textbf{Literature Review}}
A literature review was conducted to evaluate existing systems in supply chain risk management, graph-based digital twins, and agentic artificial intelligence applications.

Supply chain digital twins have emerged as a powerful tool for monitoring physical networks. By representing warehouses, ports, and lanes as nodes and edges in a network graph, digital twins can maintain a virtual representation of stock levels and transit states. Most traditional digital twins, however, are static and reactive, requiring manual data updates and predefined decision trees to calculate alternative routes.

Pathfinding algorithms, particularly Dijkstra's algorithm, are widely used in logistics to find the shortest or cheapest path between origins and destinations. In a dynamic logistics network, edge weights (representing transit time and cost) fluctuate due to disruptions. Applying Dijkstra's algorithm to dynamic graphs requires an execution engine that can propagate shocks downstream (e.g., from a chokepoint to connected ports) to calculate effective weights.

Recent advancements in Large Language Models (LLMs) and multi-agent frameworks (such as CrewAI) present opportunities to make supply chain digital twins proactive. Agentic AI systems can assign specific roles (e.g., monitoring news, optimizing routes, analyzing costs, and generating executive reports) to specialized software agents. These agents can use tools to query the digital twin's database, run mathematical optimization engines, and draft human-readable recommendations. This hybrid approach—combining the symbolic reasoning of graphs and pathfinding with the linguistic capability of LLMs—bridges the gap between complex analytical calculations and operational decision-making.

Based on the review of over 10 recent research papers, the key research gaps identified include:
\begin{itemize}
    \item Lacks automated translation of unstructured news alerts (e.g., geopolitical reports) into quantitative logistics penalties.
    \item Lacks coordination between route optimization (symbolic algorithms) and strategic inventory management (safety stock, restocking).
    \item A need for human-in-the-loop dashboard triggers that present visual recommendations alongside financial and temporal trade-offs.
\end{itemize}

\section[Motivation]{\textbf{Motivation}}
The motivation for this project stems from recent real-world maritime crises that disrupted global trade. The 2021 Suez Canal blockage by the container ship \textit{Ever Given} and the 2023–2024 Red Sea shipping crisis (resulting from attacks on commercial vessels) highlighted the fragility of global shipping corridors. These events caused container freight rates to surge and delayed deliveries to Indian ports by 10 to 15 days. For Indian SMEs, which often operate with limited safety stock, these delays quickly led to inventory stockouts, halting production lines and causing financial losses. 

Developing a system that can dynamically simulate these friction shocks and recommend optimized rerouting paths before cargo arrives at disrupted zones is crucial for maintaining supply chain resilience.

\section[Problem statement]{\textbf{Problem statement}}
The problem addressed in this project is defined as:
\begin{quote}
"To design and implement an Agentic AI-driven Digital Twin framework that dynamically models maritime logistics friction at global chokepoints, simulates the downstream propagation of transit delays and cost spikes, and autonomously recommends optimized alternative routing paths to prevent inventory stockouts in regional Indian SME markets."
\end{quote}

\section[Objectives]{\textbf{Objectives}}
The objectives of the project are:
\begin{enumerate}
\item To design a multi-tier directed network graph model (Digital Twin) in NetworkX representing the shipping lanes from international suppliers through maritime chokepoints to Indian SME markets.
\item To develop an Execution Engine that injects simulated friction events (shocks) into the graph, propagates them downstream, and runs Dijkstra's algorithm to find optimal alternate paths based on time or cost criteria.
\item To implement a multi-agent orchestration layer (using CrewAI and Google Gemini) that parses environment signals, evaluates routing options, checks market inventory levels, and suggests emergency restocking actions.
\end{enumerate}

\section[Brief Methodology of the project]{\textbf{Brief Methodology of the project}}
The methodology is executed through the following components:
\begin{itemize}
    \item \textbf{Data Ingestion Layer:} Models 50 simulated friction event scenarios and maps them to quantitative edge multipliers using predefined semantic rules.
    \item \textbf{Digital Twin Layer:} Constructs a 14-node, 19-edge directed graph in NetworkX storing baseline transit times, logistics costs, and friction penalties.
    \item \textbf{Execution Engine:} Implements shock injection logic and Dijkstra's algorithm, comparing the "Baseline Graph" and "Shocked Graph" to compute transit and cost deltas.
    \item \textbf{Agentic Intelligence:} Coordinates a sequential workflow of four agents (Friction Sentry, Logistics Optimizer, Cost Analyst, and SME Communicator) to analyze shocks, check inventory, and write summaries.
    \item \textbf{Presentation Layer:} Renders the results on an interactive React dashboard with Leaflet map overlays and real-time status alerts.
\end{itemize}

\section[Assumptions made / Constraints of the project]{\textbf{Assumptions made / Constraints of the project}}
The project is bound by the following assumptions and constraints:
\begin{itemize}
\item The digital twin topology is fixed at 14 nodes and 19 edges, representing primary corridors.
\item Transit times and logistics costs are simulated averages and do not account for daily spot freight fluctuations.
\item The multi-agent orchestration layer requires active internet connectivity and access to the Google Gemini API.
\item The inventory simulation assumes a constant demand rate at each SME market.
\end{itemize}

\section[Organization of the report]{\textbf{Organization of the report}}
This report is organized as follows:
\begin{itemize}
\item \textbf{Chapter 2} discusses the theoretical background, including graph representation, pathfinding algorithms, agentic AI frameworks, and inventory management theory.
\item \textbf{Chapter 3} detail the system design and architecture, outlining the 5-layer design, graph schemas, and the multi-agent orchestration setup.
\item \textbf{Chapter 4} presents the implementation details, including the Python scripts, API endpoints, and dashboard files.
\item \textbf{Chapter 5} provides the verification and results, documenting the test suite, simulated shock outcomes, and trade-off analyses.
\item \textbf{Chapter 6} summarizes the conclusions, details the learning outcomes, and suggests future improvements.
\end{itemize}